% !TeX program = xelatex
% UTF-8; compile twice with XeLaTeX. No external images or data files needed.
\documentclass[a4paper,landscape,10pt]{article}
\usepackage[margin=14mm,headheight=14pt,headsep=5mm,footskip=7mm]{geometry}
\usepackage{fontspec}
\setmainfont{DejaVu Sans}
\newfontfamily\ipafont{DejaVu Sans}
\newfontfamily\arabicfont[Script=Arabic]{DejaVu Sans}
\TeXXeTstate=1
\newcommand{\LR}[1]{\beginL#1\endL}
\newcommand{\AR}[1]{\leavevmode\beginR{\arabicfont\fontsize{11}{15}\selectfont #1}\endR}
\usepackage[table]{xcolor}
\usepackage{array,longtable,booktabs,fancyhdr}
\definecolor{ink}{HTML}{18384D}
\definecolor{band}{HTML}{E7EFF4}
\pagestyle{fancy}\fancyhf{}
\fancyhead[L]{\small\color{ink}Parlez Vous QSO?! · Arabic / Arabisch}
\fancyhead[R]{\small Englische Phrase · Arabische Übersetzung · IPA}
\fancyfoot[L]{\scriptsize Vorlage: Arabic v 1.0 · Credits: Cpt. Mustapha Landoulsi -- DL1BDF}
\fancyfoot[R]{\small\thepage}
\setlength{\parindent}{0pt}
\setlength{\tabcolsep}{3mm}
\renewcommand{\arraystretch}{1.25}
\newcommand{\ipa}[1]{{\ipafont[#1]}}
\begin{document}
{\LARGE\bfseries\color{ink}Parlez Vous QSO?!}\hfill{\large Arabic / Arabisch}

\vspace{2mm}
{\small QSO-Sprachführer mit hocharabischer IPA-Lautschrift. Die Lautschrift bezieht sich auf die arabische Spalte.}

\vspace{2mm}
{\footnotesize \textbf{Leseschlüssel:} \ipa{ˈ} Betonung, \ipa{ː} langer Vokal; \ipa{ʔ} Stimmabsatz,
\ipa{ʕ} ʿAin, \ipa{ħ} stimmloser Rachenlaut, \ipa{x, ɣ} wie ch beziehungsweise dessen stimmhaftes Gegenstück;
\ipa{sˤ, dˤ, tˤ, ðˤ} emphatische Laute. Arabisch rechts nach links; IPA und lateinische Kürzel links nach rechts.}

\vspace{1mm}
\fontsize{9}{11.8}\selectfont
\begin{longtable}{@{}>{\raggedright\arraybackslash}p{13mm}>{\raggedright\arraybackslash}p{73mm}>{\raggedleft\arraybackslash}p{75mm}>{\raggedright\arraybackslash}p{90mm}@{}}
\toprule
\textbf{Ref.} & \textbf{English phrase} & \textbf{Arabisch} & \textbf{Lautschrift (IPA)}\\
\midrule
\endfirsthead
\toprule
\textbf{Ref.} & \textbf{English phrase} & \textbf{Arabisch} & \textbf{Lautschrift (IPA)}\\
\midrule
\endhead
\midrule
\multicolumn{4}{r}{\scriptsize Fortsetzung auf der nächsten Seite}\\
\endfoot
\bottomrule
\endlastfoot
\rowcolor{band} 1.0 & \textbf{Opening Greetings} & \multicolumn{2}{r}{\AR{بدء الاتصال}} \\*
1.1 & I am very pleased to meet you. & \AR{يسعدني أن أتواصل معك.} & \ipa{ˈjusʕiduni ʔan ʔataˈwaːsˤal ˈmaʕak} \\
1.2 & Thank you for your call & \AR{شكرا لك على ندائك.} & \ipa{ˈʃukran lak ʕalaː niˈdaːʔik} \\
1.3 & Good morning & \AR{صباح الخير.} & \ipa{sˤaˈbaːħ alˈxajr} \\
1.31 & Good day & \AR{نهارك سعيد.} & \ipa{naˈhaːrak saˈʕiːd} \\
1.32 & Good afternoon & \AR{طاب عصرُك.} & \ipa{tˤaːb ˈʕasˤruk} \\
1.33 & Good evening & \AR{مساء الخير.} & \ipa{maˈsaːʔ alˈxajr} \\
1.4 & I know only a little Arabic but I enjoy QSOs in Arabic. & \AR{أتحدث العربية قليلا، لكنني أستمتع بإجراء الاتصالات بالعربية.} & \ipa{ʔataˈħaddaθ alʕaraˈbijja qaˈliːlan ˈlaːkinnani ʔasˈtamtiʕ biʔidʒˈraːʔ alittiˈsˤaːlaːt bilʕaraˈbijja} \\
\rowcolor{band} 2.0 & \textbf{Operator} & \multicolumn{2}{r}{\AR{المشغّل}} \\*
2.1 & My name is Peter. & \AR{اسمي بيتر.} & \ipa{ˈʔismi ˈbiːtar} \\
2.2 & I have been licensed since 1960. & \AR{أنا مرخّص منذ عام \LR{1960}.} & \ipa{ˈʔanaː muˈraxxasˤ munð ʕaːm ʔalf waˈtisʕ miˈʔa wa sitˈtiːn} \\
2.3 & I have been licensed for 43 years. & \AR{أنا مرخّص منذ \LR{43} سنة.} & \ipa{ˈʔanaː muˈraxxasˤ munð θaˈlaːθ waʔarbaˈʕiːn ˈsana} \\
2.4 & I live in Sittingbourne. & \AR{أسكن في سيتنغبورن.} & \ipa{ˈʔaskun fiː ˈsitiŋbuːrn} \\
2.5 & Sittingbourne is a village. & \AR{سيتنغبورن قرية.} & \ipa{ˈsitiŋbuːrn ˈqarja} \\
2.51 & Sittingbourne is a town. & \AR{سيتنغبورن بلدة.} & \ipa{ˈsitiŋbuːrn ˈbalda} \\
2.52 & Sittingbourne is a city. & \AR{سيتنغبورن مدينة.} & \ipa{ˈsitiŋbuːrn maˈdiːna} \\
2.6 & Sittingbourne is between London and Dover. & \AR{تقع سيتنغبورن بين لندن ودوفر.} & \ipa{ˈtaqaʕ ˈsitiŋbuːrn bajn ˈlandan waˈduːfar} \\
2.7 & I live in the South of England. & \AR{أسكن في جنوب إنجلترا.} & \ipa{ˈʔaskun fiː dʒaˈnuːb ʔinˈdʒiltiraː} \\
2.71 & I live in the North of England. & \AR{أسكن في شمال إنجلترا.} & \ipa{ˈʔaskun fiː ʃaˈmaːl ʔinˈdʒiltiraː} \\
2.72 & I live in the East of England. & \AR{أسكن في شرق إنجلترا.} & \ipa{ˈʔaskun fiː ʃarq ʔinˈdʒiltiraː} \\
2.73 & I live in West of England. & \AR{أسكن في غرب إنجلترا.} & \ipa{ˈʔaskun fiː ɣarb ʔinˈdʒiltiraː} \\
2.74 & I live in South East England. & \AR{أسكن في جنوب شرق إنجلترا.} & \ipa{ˈʔaskun fiː dʒaˈnuːb ʃarq ʔinˈdʒiltiraː} \\
2.75 & I live in South West England. & \AR{أسكن في جنوب غرب إنجلترا.} & \ipa{ˈʔaskun fiː dʒaˈnuːb ɣarb ʔinˈdʒiltiraː} \\
2.76 & I live in North East England. & \AR{أسكن في شمال شرق إنجلترا.} & \ipa{ˈʔaskun fiː ʃaˈmaːl ʃarq ʔinˈdʒiltiraː} \\
2.77 & I live in North West England. & \AR{أسكن في شمال غرب إنجلترا.} & \ipa{ˈʔaskun fiː ʃaˈmaːl ɣarb ʔinˈdʒiltiraː} \\
2.78 & I live in central England. & \AR{أسكن في وسط إنجلترا.} & \ipa{ˈʔaskun fiː ˈwasatˤ ʔinˈdʒiltiraː} \\
2.8 & I live in the country. & \AR{أسكن في الريف.} & \ipa{ˈʔaskun firˈriːf} \\
2.81 & I live by the sea. & \AR{أسكن قرب البحر.} & \ipa{ˈʔaskun qurb alˈbaħr} \\
2.82 & I live near mountains. & \AR{أسكن قرب الجبال.} & \ipa{ˈʔaskun qurb aldʒiˈbaːl} \\
2.9 & Sittingbourne is 50Km from London. & \AR{تبعد سيتنغبورن \LR{50} كيلومترا عن لندن.} & \ipa{ˈtabʕud ˈsitiŋbuːrn xamˈsiːn kiːluˈmitran ʕan ˈlandan} \\
2.11 & I am retired. & \AR{أنا متقاعد.} & \ipa{ˈʔanaː mutaˈqaːʕid} \\
2.12 & My home call sign is M0ECS. & \AR{علامة ندائي في بلدي هي \LR{M0ECS}.} & \ipa{ʕaˈlaːmat niˈdaːʔi fiː ˈbaladi hiˈja em ˈsˤifr iː siː es} \\
\rowcolor{band} 3.0 & \textbf{Station} & \multicolumn{2}{r}{\AR{المحطة}} \\*
3.1 & My radio is an Icom 706. & \AR{جهاز اللاسلكي لديّ هو \LR{Icom} \LR{706}.} & \ipa{dʒiˈhaːz allaːˈsilki laˈdajja ˈhuwa ˈaɪkɒm sabʕ miˈʔa waˈsitta} \\
3.2 & My station is all homebrew (homemade). & \AR{صنعت جميع أجهزة محطتي بنفسي.} & \ipa{sˤaˈnaʕt dʒaˈmiːʕ ˈʔadʒhizat maˈħatˤtˤati biˈnafsi} \\
3.3 & I am running 5 Watts. & \AR{أرسل بقدرة خمسة واط.} & \ipa{ˈʔursil biˈqudrat ˈxamsat waːtˤ} \\
3.4 & My antenna is a long wire. & \AR{الهوائي لديّ سلك طويل.} & \ipa{alhawaːˈʔi laˈdajja silk tˤaˈwiːl} \\
3.41 & My antenna is a dipole. & \AR{الهوائي لديّ ثنائي القطب.} & \ipa{alhawaːˈʔi laˈdajja θuˈnaːʔi alˈqutˤb} \\
3.42 & My antenna is a G5RV. & \AR{الهوائي لديّ من نوع \LR{G5RV}.} & \ipa{alhawaːˈʔi laˈdajja min nawʕ dʒiː ˈxamsa ɑːr viː} \\
3.43 & My antenna is a doublet. & \AR{الهوائي لديّ ثنائي القطب من نوع \LR{Doublet}.} & \ipa{alhawaːˈʔi laˈdajja θuˈnaːʔi alˈqutˤb min nawʕ ˈdʌblɪt} \\
3.44 & My antenna is a delta loop. & \AR{الهوائي لديّ حلقي مثلّث.} & \ipa{alhawaːˈʔi laˈdajja ˈħalqi muˈθallaθ} \\
3.45 & My antenna is a vertical. & \AR{الهوائي لديّ عمودي.} & \ipa{alhawaːˈʔi laˈdajja ʕaˈmuːdi} \\
3.46 & My antenna is a log periodic. & \AR{الهوائي لديّ لوغاريتمي دوري.} & \ipa{alhawaːˈʔi laˈdajja luɣaːˈriːtmi ˈdawri} \\
3.47 & My antenna is a Yagi beam. & \AR{الهوائي لديّ ياغي اتجاهي.} & \ipa{alhawaːˈʔi laˈdajja ˈjaːɣi ittiˈdʒaːhi} \\
3.48 & My antenna is a 5 element beam. & \AR{الهوائي لديّ اتجاهي ذو خمسة عناصر.} & \ipa{alhawaːˈʔi laˈdajja ittiˈdʒaːhi ðuː ˈxamsat ʕaˈnaːsˤir} \\
\rowcolor{band} 4.0 & \textbf{Conditions} & \multicolumn{2}{r}{\AR{ظروف الانتشار}} \\*
4.1 & Thank you for my 599 report. & \AR{شكرا على تقرير الإشارة \LR{599}.} & \ipa{ˈʃukran ʕalaː taqˈriːr alʔiˈʃaːra ˈxamsa ˈtisʕa ˈtisʕa} \\
4.2 & The band conditions are good. & \AR{ظروف الانتشار على النطاق جيدة.} & \ipa{ðˤuˈruːf alintiˈʃaːr ʕalanˈnitˤaːq ˈdʒajjida} \\
4.3 & The band conditions are bad. & \AR{ظروف الانتشار على النطاق سيئة.} & \ipa{ðˤuˈruːf alintiˈʃaːr ʕalanˈnitˤaːq ˈsajjiʔa} \\
4.4 & There is strong interference here. & \AR{يوجد تداخل قوي هنا.} & \ipa{ˈjuːdʒad taˈdaːxul qaˈwij ˈhunaː} \\
4.5 & Can you change to the 80m band? & \AR{هل يمكنك الانتقال إلى نطاق \LR{80} مترا؟} & \ipa{hal ˈjumkinuk alintiˈqaːl ʔilaː niˈtˤaːq θamaˈniːn ˈmitran} \\
4.6 & Please change frequency to 3.558 MHz. & \AR{يرجى تغيير التردد إلى \LR{3.558} \LR{MHz}.} & \ipa{ˈjurdʒaː taɣˈjiːr attaˈraddud ʔilaː θaˈlaːθa ˈfaːsˤila ˈxamsa ˈxamsa θaˈmaːnija ˈmiːɣaˌhirts} \\
\rowcolor{band} 5.0 & \textbf{General Queries and Comments} & \multicolumn{2}{r}{\AR{أسئلة وتعليقات عامة}} \\*
5.1 & What does <\#\#\#\#\#> mean please? & \AR{ما معنى \LR{<\#\#\#\#\#>}، من فضلك؟} & \ipa{maː ˈmaʕnaː … min ˈfadˤlik} \\
5.2 & Please wait a moment, I need to check my Arabic dictionary. & \AR{انتظر لحظة من فضلك، سأراجع القاموس العربي.} & \ipa{ˈʔintaðˤir ˈlaħðˤa min ˈfadˤlik saʔuˈraːdʒiʕ alqaːˈmuːs alʕaˈrabi} \\
5.3 & How do you say, <\#\#\#\#\#> in Arabic ? & \AR{كيف تقول \LR{<\#\#\#\#\#>} بالعربية؟} & \ipa{kajf taˈquːl … bilʕaraˈbijja} \\
5.5 & My favourite modes of operation are CW / RTTY / PSK / MT63 / Olivia / Hellschreiber / Digital Modes / and SSB & \AR{أنماط التشغيل المفضلة لديّ هي: \LR{CW} / \LR{RTTY} / \LR{PSK} / \LR{MT63} / \LR{Olivia} / \LR{Hellschreiber} / الأنماط الرقمية / \LR{SSB}.} & \ipa{ʔanˈmaːtˤ attaʃˈɣiːl almufadˤˈdˤala laˈdajja hiˈja siː ˈdʌbljuː / ɑːr tiː tiː waɪ / piː es keɪ / em tiː ˈsitta wa sitˈtiːn / oˈlivia / ˈhɛlʃraɪbɐ / alʔanˈmaːtˤ arraˈqamijja / es es biː} \\
5.6 & My best DX is Australia. & \AR{أبعد اتصال لي كان مع أستراليا.} & \ipa{ˈʔabʕad ittiˈsˤaːl liː kaːn maʕ ʔusˈtraːljaː} \\
5.7 & I enjoy working contests. & \AR{أستمتع بالمشاركة في المسابقات.} & \ipa{ʔasˈtamtiʕ bilmuˈʃaːraka fiː almusaːbaˈqaːt} \\
5.8 & I like working QRP. & \AR{أحب العمل بقدرة منخفضة، أي \LR{QRP}.} & \ipa{ʔuˈħibb alˈʕamal biˈqudra munˈxafidˤa ʔaj kjuː ɑːr piː} \\
\rowcolor{band} 6.0 & \textbf{The Weather} & \multicolumn{2}{r}{\AR{الطقس}} \\*
6.1 & It is sunny. & \AR{الجو مشمس.} & \ipa{alˈdʒaww ˈmuʃmis} \\
6.2 & It is windy. & \AR{الرياح قوية.} & \ipa{arˈrijaːħ qaˈwijja} \\
6.3 & It is raining. & \AR{تمطر الآن.} & \ipa{ˈtumtˤir alˈʔaːn} \\
6.4 & It is snowing. & \AR{يتساقط الثلج.} & \ipa{jataˈsaːqatˤ aθˈθaldʒ} \\
6.5 & It is very hot. & \AR{الجو حار جدا.} & \ipa{alˈdʒaww ħaːr ˈdʒiddan} \\
6.6 & It is very cold. & \AR{الجو بارد جدا.} & \ipa{alˈdʒaww ˈbaːrid ˈdʒiddan} \\
6.7 & The temperature is 8C. & \AR{درجة الحرارة ثماني درجات مئوية.} & \ipa{ˈdaradʒat alħaˈraːra θaˈmaːni daraˈdʒaːt miʔaˈwijja} \\
6.8 & It is overcast. & \AR{السماء ملبّدة بالغيوم.} & \ipa{assaˈmaːʔ muˈlabbada bilɣuˈjuːm} \\
6.9 & It is foggy. & \AR{يوجد ضباب.} & \ipa{ˈjuːdʒad dˤaˈbaːb} \\
6.11 & It has been raining all day. & \AR{استمر المطر طوال اليوم.} & \ipa{ʔistaˈmarr alˈmatˤar tˤaˈwaːl alˈjawm} \\
6.12 & There is 20cm of snow here. & \AR{يوجد هنا ثلج بسمك \LR{20} سنتيمترا.} & \ipa{ˈjuːdʒad ˈhunaː θaldʒ biˈsumk ʕiʃˈriːn santiˈmitran} \\
\rowcolor{band} 8.0 & \textbf{QSL} & \multicolumn{2}{r}{\AR{بطاقات \LR{QSL}}} \\*
8.1 & My QSL card is sure via the bureau & \AR{سأرسل لك بطاقة \LR{QSL} بالتأكيد عبر المكتب.} & \ipa{saˈʔursil lak biˈtˤaːqat kjuː es el bittaʔˈkiːd ʕabr alˈmaktab} \\
8.2 & QSL via eQSL & \AR{أرسل بطاقة \LR{QSL} عبر \LR{eQSL}.} & \ipa{ˈʔursil biˈtˤaːqat kjuː es el ʕabr iː kjuː es el} \\
8.3 & No eQSL please & \AR{لا ترسل \LR{eQSL} من فضلك.} & \ipa{laː ˈtursil iː kjuː es el min ˈfadˤlik} \\
8.4 & I will send you my card direct. & \AR{سأرسل لك بطاقتي مباشرة.} & \ipa{saˈʔursil lak biˈtˤaːqati muˈbaːʃara} \\
8.5 & What is your address? & \AR{ما عنوانك؟} & \ipa{maː ʕunˈwaːnuk} \\
8.6 & Please send your QSL card direct. & \AR{أرسل بطاقة \LR{QSL} مباشرة من فضلك.} & \ipa{ˈʔarsil biˈtˤaːqat kjuː es el muˈbaːʃara min ˈfadˤlik} \\
8.7 & My address is... & \AR{عنواني هو…} & \ipa{ʕunˈwaːni ˈhuwa …} \\
8.8 & My address is shown on QSL.com / EHam.net & \AR{عنواني موجود على \LR{QSL.com} / \LR{EHam.net}.} & \ipa{ʕunˈwaːni mawˈdʒuːd ʕalaː kjuː es el nuqˈtˤa kɒm / iː hæm nuqˈtˤa net} \\
8.9 & Sorry I do not send QSL cards. & \AR{آسف، لا أرسل بطاقات \LR{QSL}.} & \ipa{ˈʔaːsif laː ˈʔursil bitˤaːˈqaːt kjuː es el} \\
8.11 & Please send QSL card via G1JQH. & \AR{أرسل بطاقة \LR{QSL} عبر \LR{G1JQH} من فضلك.} & \ipa{ˈʔarsil biˈtˤaːqat kjuː es el ʕabr dʒiː ˈwaːħid dʒeɪ kjuː eɪtʃ min ˈfadˤlik} \\
8.12 & My WACRAL / GQRP / DXCC / FISTS number is 1234. & \AR{رقمي في \LR{WACRAL} / \LR{GQRP} / \LR{DXCC} / \LR{FISTS} هو \LR{1234}.} & \ipa{ˈraqami fiː ˈwakral / dʒiː kjuː ɑːr piː / diː eks siː siː / fɪsts ˈhuwa ʔalf wa miʔaˈtaːn waˈʔarbaʕ waθalaˈθuːn} \\
\rowcolor{band} 9.0 & \textbf{Final Greetings} & \multicolumn{2}{r}{\AR{ختام الاتصال}} \\*
9.1 & Thank you. & \AR{شكرا.} & \ipa{ˈʃukran} \\
9.2 & Thank you for a very enjoyable contact. & \AR{شكرا على هذا الاتصال الممتع.} & \ipa{ˈʃukran ʕalaː ˈhaːðaː alittiˈsˤaːl alˈmumtiʕ} \\
9.3 & I enjoyed our chat. & \AR{استمتعت بحديثنا.} & \ipa{ʔistamˈtaʕt biħaˈdiːθinaː} \\
9.4 & I enjoyed our QSO in Arabic. & \AR{استمتعت باتصالنا بالعربية.} & \ipa{ʔistamˈtaʕt biittiˈsˤaːlinaː bilʕaraˈbijja} \\
9.5 & It has been great fun. & \AR{كان ذلك ممتعا جدا.} & \ipa{kaːn ˈðaːlik ˈmumtiʕan ˈdʒiddan} \\
9.6 & If you would like a regular contact with me, you can contact me via email. & \AR{إذا أردت التواصل معي بانتظام، يمكنك مراسلتي بالبريد الإلكتروني.} & \ipa{ʔiˈðaː ʔaˈradt attaˈwaːsˤul ˈmaʕiː bintiˈðˤaːm ˈjumkinuk muraːˈsalati bilbaˈriːd alʔilikˈtroːni} \\
9.7 & My email address is... & \AR{عنوان بريدي الإلكتروني هو…} & \ipa{ʕunˈwaːn baˈriːdi alʔilikˈtroːni ˈhuwa …} \\
9.8 & My email address is shown on QSL.com / EHam.net & \AR{عنوان بريدي الإلكتروني موجود على \LR{QSL.com} / \LR{EHam.net}.} & \ipa{ʕunˈwaːn baˈriːdi alʔilikˈtroːni mawˈdʒuːd ʕalaː kjuː es el nuqˈtˤa kɒm / iː hæm nuqˈtˤa net} \\
9.9 & I look forward to meeting you again. & \AR{أتطلع إلى لقائك مرة أخرى.} & \ipa{ʔataˈtˤallaʕ ʔilaː liˈqaːʔik ˈmarra ˈʔuxraː} \\
9.11 & My best wishes to you and your family. & \AR{أطيب تمنياتي لك ولأسرتك.} & \ipa{ˈʔatˤjab tamanˈnjaːti lak waliˈʔusratik} \\
9.12 & Best wishes to your wife. & \AR{أطيب تمنياتي لزوجتك.} & \ipa{ˈʔatˤjab tamanˈnjaːti liˈzawdʒatik} \\
9.121 & Best wishes to your husband. & \AR{أطيب تمنياتي لزوجك.} & \ipa{ˈʔatˤjab tamanˈnjaːti liˈzawdʒik} \\
9.13 & Good luck. & \AR{حظا سعيدا.} & \ipa{ˈħaðˤðˤan saˈʕiːdan} \\
9.14 & Good health. & \AR{أتمنى لك دوام الصحة.} & \ipa{ʔataˈmannaː lak daˈwaːm asˈsˤiħħa} \\
9.15 & Goodbye & \AR{إلى اللقاء.} & \ipa{ʔilaː alliˈqaːʔ} \\
9.16 & I wish you an enjoyable afternoon. & \AR{أتمنى لك فترة ما بعد الظهر ممتعة.} & \ipa{ʔataˈmannaː lak ˈfatrat maː baʕd aðˤˈðˤuhr ˈmumtiʕa} \\
9.161 & I wish you an enjoyable evening. & \AR{أتمنى لك مساء ممتعا.} & \ipa{ʔataˈmannaː lak maˈsaːʔan ˈmumtiʕan} \\
9.162 & I wish you an enjoyable weekend. & \AR{أتمنى لك عطلة نهاية أسبوع ممتعة.} & \ipa{ʔataˈmannaː lak ˈʕutˤlat niˈhaːjat ʔusˈbuːʕ ˈmumtiʕa} \\
9.17 & Enjoy your holiday. & \AR{استمتع بإجازتك.} & \ipa{ʔisˈtamtiʕ biʔiˈdʒaːzatik} \\
9.18 & Happy New Year. & \AR{سنة جديدة سعيدة!} & \ipa{ˈsana dʒaˈdiːda saˈʕiːda} \\
\rowcolor{band} 10.0 & \textbf{Country} & \multicolumn{2}{r}{\AR{البلد}} \\*
10.1 & Albania & \AR{ألبانيا} & \ipa{ʔalˈbaːnjaː} \\
10.2 & Austria & \AR{النمسا} & \ipa{anˈnimsaː} \\
10.3 & Belgium & \AR{بلجيكا} & \ipa{balˈdʒiːkaː} \\
10.4 & Belorussia & \AR{بيلاروس} & \ipa{biːlaˈruːs} \\
10.5 & Bosnia \& Herzegovina & \AR{البوسنة والهرسك} & \ipa{alˈbuːsna walˈhirsik} \\
10.6 & Bulgaria & \AR{بلغاريا} & \ipa{bulˈɣaːrjaː} \\
10.7 & Croatia & \AR{كرواتيا} & \ipa{kuruˈwaːtjaː} \\
10.8 & Czech Republic & \AR{جمهورية التشيك} & \ipa{dʒumhuˈrijjat atˈtʃiːk} \\
10.9 & Denmark & \AR{الدنمارك} & \ipa{adˈdinmaːrk} \\
10.11 & Eire & \AR{أيرلندا} & \ipa{ʔajrˈlandaː} \\
10.12 & England & \AR{إنجلترا} & \ipa{ʔinˈdʒiltiraː} \\
10.13 & Estonia & \AR{إستونيا} & \ipa{ʔisˈtuːnjaː} \\
10.14 & Finland & \AR{فنلندا} & \ipa{finˈlandaː} \\
10.15 & France & \AR{فرنسا} & \ipa{faˈransaː} \\
10.16 & Germany & \AR{ألمانيا} & \ipa{ʔalˈmaːnjaː} \\
10.17 & Gibraltar & \AR{جبل طارق} & \ipa{ˈdʒabal ˈtˤaːriq} \\
10.18 & Greece & \AR{اليونان} & \ipa{aljuːˈnaːn} \\
10.19 & Holland & \AR{هولندا} & \ipa{huːˈlandaː} \\
10.21 & Hungary & \AR{المجر} & \ipa{alˈmadʒar} \\
10.22 & Iceland & \AR{آيسلندا} & \ipa{ʔajsˈlandaː} \\
10.23 & Northern Ireland & \AR{أيرلندا الشمالية} & \ipa{ʔajrˈlandaː aʃʃamaːˈlijja} \\
10.24 & Isle of Man & \AR{جزيرة مان} & \ipa{dʒaˈziːrat maːn} \\
10.25 & Italy & \AR{إيطاليا} & \ipa{ʔiːˈtˤaːljaː} \\
10.26 & Latvia & \AR{لاتفيا} & \ipa{ˈlaːtfjaː} \\
10.27 & Liechtenstein & \AR{ليختنشتاين} & \ipa{liːxtinˈʃtaːjn} \\
10.28 & Lithuania & \AR{ليتوانيا} & \ipa{liːtuˈwaːnjaː} \\
10.29 & Luxembourg & \AR{لوكسمبورغ} & \ipa{luːksamˈbuːrɣ} \\
10.31 & Macedonia & \AR{مقدونيا} & \ipa{maqˈduːnjaː} \\
10.32 & Moldavia & \AR{مولدوفا} & \ipa{muːlˈduːfaː} \\
10.33 & Netherlands & \AR{هولندا} & \ipa{huːˈlandaː} \\
10.34 & Norway & \AR{النرويج} & \ipa{annarˈwiːdʒ} \\
10.35 & Poland & \AR{بولندا} & \ipa{buːˈlandaː} \\
10.36 & Portugal & \AR{البرتغال} & \ipa{alburtuˈɣaːl} \\
10.37 & Romania & \AR{رومانيا} & \ipa{ruːˈmaːnjaː} \\
10.38 & Russian Federation & \AR{الاتحاد الروسي} & \ipa{alittiˈħaːd arˈruːsi} \\
10.39 & Scotland & \AR{إسكتلندا} & \ipa{ʔiskutˈlandaː} \\
10.41 & Slovakia & \AR{سلوفاكيا} & \ipa{sluːˈfaːkjaː} \\
10.42 & Slovenia & \AR{سلوفينيا} & \ipa{sluːˈfiːnjaː} \\
10.43 & Spain & \AR{إسبانيا} & \ipa{ʔisˈbaːnjaː} \\
10.44 & Svalbard & \AR{سفالبارد} & \ipa{sfaːlˈbaːrd} \\
10.45 & Sweden & \AR{السويد} & \ipa{assuˈwajd} \\
10.46 & Switzerland & \AR{سويسرا} & \ipa{suˈwajsraː} \\
10.47 & Turkey & \AR{تركيا} & \ipa{ˈturkjaː} \\
10.48 & Ukraine & \AR{أوكرانيا} & \ipa{ʔuːˈkraːnjaː} \\
10.49 & Wales & \AR{ويلز} & \ipa{wajlz} \\
\end{longtable}
\normalsize
\clearpage
\section*{Hinweise zur Überarbeitung und Aussprache}
\fontsize{8.5}{10}\selectfont\raggedright
\textbf{Grundlage.} „Arabic v 1.0.doc“, Credit: Cpt. Mustapha Landoulsi, DL1BDF.
Alle 142 englischen Phrasen und Ländernamen sind enthalten. Nummernsprünge und der fehlende Abschnitt 7
bleiben erhalten; 132 wurde zu 1.32. Die arabische Spalte ist redaktionell überarbeitet und ergänzt.
Die alte lateinische Umschrift wurde durch IPA ersetzt, da sie stellenweise nicht zum arabischen Text passte.

\smallskip
\textbf{Schrift und Leserichtung.} Arabische Wörter werden mit verbundenen Buchstaben von rechts nach links
gesetzt und rechtsbündig ausgerichtet. Zahlen, Modellnamen, Rufzeichen, Webadressen und lateinische Kürzel
behalten innerhalb dieser Spalte ihre Richtung von links nach rechts.
Die gesamte Vergleichstabelle bleibt in der bisherigen Spaltenfolge Ref. / Englisch / Arabisch / IPA.
Der Quelltext enthält arabische Unicode-Zeichen in logischer Schreibreihenfolge, nicht manuell umgedreht.

\smallskip
\textbf{Sprachform.} Grundlage der neuen Spalte ist allgemein verständliches modernes Hocharabisch.
Die Vorlage mischte Dialektformen wie eisch, misch, kuwayyis und mabsut mit Hocharabisch.
Die neue Fassung vereinheitlicht Ausdruck und Grammatik; die IPA-Spalte beschreibt diese neue Fassung.
Sie ist keine Wiedergabe eines bestimmten regionalen Dialekts.

\smallskip
\textbf{Wichtige Berichtigungen.} Ahmed wurde gemäß der englischen Spalte zu Peter (2.1);
Sittingbourne wurde in 2.4 ergänzt. Ost/West waren im arabischen Text als Nord/Ost wiedergegeben (2.72/2.73).
Die fehlerhaften und doppelten Mischtexte in 2.9 und 3.2 wurden neu formuliert.
4.1 enthält jetzt tatsächlich 599; 4.6 enthält 3.558 MHz statt der ausgeschriebenen falschen 3.553.
5.3 fragt nun, wie man etwas auf Arabisch sagt, statt nur nach seiner Bedeutung.
Alle Betriebsarten in 5.5 sind enthalten, einschließlich MT63.
6.3 bedeutet jetzt Regen, 6.8 bedeckten Himmel, 6.11 Regen während des ganzen Tages.
Unvollständige Temperaturformulierungen in 6.5/6.6 wurden bereinigt.

\smallskip
\textbf{Fehlende Angaben und Abschiedsformeln.} Die Überschriften 4.0, 5.0, 8.0 und 10.0 sowie der Eintrag
8.12 wurden ergänzt. Rufzeichen und Webadressen wurden in den arabischen Sätzen vervollständigt.
9.12 bezieht sich nun auf die Ehefrau statt die ganze Familie; 9.121 auf den Ehemann.
9.162 wünscht ausdrücklich ein schönes Wochenende statt allgemein eine freie Zeit.
Die Ausdrücke für Nachmittag und Abend sind getrennt: 1.32 verwendet einen möglichen Nachmittagswunsch
(„Möge dein Nachmittag angenehm sein“); 9.16 einen ausdrücklich auf die Zeit nach Mittag bezogenen Wunsch.
Es gibt mehrere regionale Grußkonventionen; die Behauptung, ein Nachmittagswunsch sei auf Arabisch unmöglich,
wurde entfernt. 9.18 wünscht ausdrücklich ein glückliches neues Jahr.

\smallskip
\textbf{Antennenbegriffe.} Die Hinweise „kein arabischer Ausdruck“ wurden durch sinnvolle Bezeichnungen ersetzt:
langer Draht, Dipol, dreieckige Schleife, vertikales, logarithmisch-periodisches und gerichtetes Yagi-Antennensystem.
Doublet wird als Dipol mit beibehaltener englischer Typbezeichnung beschrieben; G5RV bleibt ein Modellname.
Ein Doublet kann je nach Ausführung, Länge und Speisung unterschiedliche Eigenschaften haben;
die Beschreibung legt keine bestimmte Resonanzlänge fest. Betriebsartenkürzel und Eigennamen bleiben unverändert.

\smallskip
\textbf{Ländernamen.} Irland und Deutschland fehlten und wurden ergänzt.
Griechenland stand fälschlich als „die Augen“, Portugal als „die Orangen“; beides ist berichtigt.
Nordirland, Isle of Man, Lettland/Litauen, Norwegen, Slowenien, Svalbard und Wales wurden korrigiert oder präzisiert.
Wales ist nicht Schottland. Svalbard wird als Eigenname wiedergegeben, ohne den Kommentar „kenne dieses Land nicht“;
es ist eine norwegische Inselgruppe. Belarus, Moldau und Russische Föderation verwenden gebräuchliche moderne
arabische Namen. „Macedonia“ bleibt als historischer englischer Eintrag unverändert.

\smallskip
\textbf{Aussprache und Anrede.} Die breite IPA-Hilfe markiert Betonung, Vokallänge, Verdopplung und wichtige
arabische Konsonanten. \ipa{ʕ, ħ} sind Rachenlaute, \ipa{q} ein weiter hinten artikulierter k-artiger Laut;
\ipa{ˤ} kennzeichnet Emphase. \ipa{dʒ} ist hier die Lesart von ج; regionale \ipa{ʒ} oder \ipa{ɡ} sind möglich.
Der Artikel al- ist vor Sonnenbuchstaben assimiliert, etwa in ar-riyāḥ.
Kurze Kasus- und Modusendungen werden überwiegend weggelassen; feste Formen wie shukran bleiben erhalten.
Die Anrede folgt meist der männlichen Einzahl; für eine weibliche Gesprächspartnerin ändern sich insbesondere
Suffixe und Imperative. Satzintonation, vollständige Vokalfärbung und regionale Varianten sind nicht kodiert.

\smallskip
\textbf{Namen und Zahlen.} Arabische Orts- und Ländernamen werden nach der hier gewählten arabischen Form gelesen.
Icom, Doublet, Olivia, Hellschreiber, WACRAL und FISTS verwenden fremdsprachige Lesarten.
Rufzeichen und Kürzel verwenden englische Buchstabennamen; Zahlen darin werden arabisch gelesen.
599 wird ziffernweise gelesen; die Dezimalstellen von 3.558 ebenfalls.
1234 wird als tausend plus zweihundert plus vierunddreißig gesprochen.
Offene Angaben und Platzhalter erscheinen als \ipa{…}. MHz wird mit der arabischen Lesart von Megahertz ausgesprochen.

\smallskip
\textbf{Fachhinweise.} Robin Thelwall und M. Akram Sa'adeddin: „Arabic“, Journal of the International
Phonetic Association 20(2), 1990, S. 37--41, DOI 10.1017/S0025100300004266.
Arabische Antennenbezeichnungen wurden mit dem arabischen
ITU-R-Bericht SM.2303-1 (2015) und einem Herstellertext von Tesswave zur logarithmisch-periodischen Antenne abgeglichen. Die Neufassung ist eine redaktionelle Sprach- und Aussprachehilfe.
\end{document}
